% !TeX root = ../main.tex
%======================================================================
% commands.tex - Eigene Befehle und Abkuerzungen (KORRIGIERT)
%======================================================================

% --- Mathematische Operatoren ---
\DeclareMathOperator{\diag}{diag}
\DeclareMathOperator{\re}{Re}
\DeclareMathOperator{\im}{Im}

% --- Haeufig verwendete Symbole ---
\newcommand{\R}{\mathbb{R}}
\newcommand{\C}{\mathbb{C}}
\newcommand{\jj}{\mathrm{j}}
\newcommand{\ee}{\mathrm{e}}

% --- Vektoren und Matrizen ---
\newcommand{\vect}[1]{\boldsymbol{#1}}
\newcommand{\matr}[1]{\mathbf{#1}}
\newcommand{\tens}[1]{\mathcal{#1}}

% --- Spezielle Matrizen ---
\newcommand{\ZB}{Z_B}
\newcommand{\Ydd}{Y_{dd}}
\newcommand{\Yds}{Y_{ds}}

% --- Hadamard-Operationen ---
\newcommand{\had}{\odot}
\newcommand{\haddiv}{\oslash}
\newcommand{\eleminv}{{\circ(-1)}}

% --- Abkuerzungen ---
\newcommand{\pu}{\,\text{p.u.}}
\newcommand{\spec}{\mathrm{spec}}
\newcommand{\conv}{\mathrm{konv}}

% --- Per-Unit-Groessen ---
\newcommand{\Vspec}{|V^{\spec}|}

% --- Dot-Notation ---
\newcommand{\dotV}{\dot{V}}
\newcommand{\dotS}{\dot{S}}
\newcommand{\dotM}{\dot{M}}
\newcommand{\dotW}{\dot{W}}

% --- Boxed Equation ---
\newcommand{\boxedeq}[1]{\begin{equation}\boxed{#1}\end{equation}}

% --- Theorem-Umgebungen (mit amsthm) ---
\usepackage{amsthm}
\newtheorem{theorem}{Satz}[chapter]
\newtheorem{proposition}[theorem]{Proposition}
\newtheorem{lemma}[theorem]{Lemma}
\theoremstyle{definition}
\newtheorem{definition}[theorem]{Definition}
\theoremstyle{remark}
\newtheorem{remark}[theorem]{Bemerkung}